\section{Visualization: exact midpoint-square justification}
\label{annex:p19-rat133-midpoint-square}

\begin{annexitem}{annex:p19.rat133-square-area}{thm:p19-rat133-square-area}{The source's geometric example under exact hypotheses}
\begin{proof}
Place the outer square at \((0,0),(\ell,0),(\ell,\ell),(0,\ell)\).
The consecutive side midpoints are
\[
 (\ell/2,0),\quad(\ell,\ell/2),\quad
 (\ell/2,\ell),\quad(0,\ell/2).
\]
Joining them leaves one triangle at each corner, each a right triangle
with two legs \(\ell/2\). Their areas are each
\(\frac12(\ell/2)(\ell/2)=\ell^2/8\). Their interiors are disjoint,
and together with the inner quadrilateral they partition the outer square.
The remaining area is therefore
\[
             \ell^2-4(\ell^2/8)=\ell^2/2.
\]
To verify that the quadrilateral really is a square, each side joins
midpoints of two adjacent outer sides and is the hypotenuse of one of
these congruent right-isosceles triangles. Hence its four sides have equal
length. Each corner triangle has its two acute angles equal to
\(45^\circ\); at a side midpoint, the inner angle is
\(180^\circ-45^\circ-45^\circ=90^\circ\). Thus the inner
quadrilateral has four equal sides and four right angles, and is a square.

This uses Euclidean square geometry and exact midpoints, not a measured
pixel ratio. The four triangles can also be reflected across their
hypotenuses into the inner square, whose center partitions it into four
congruent triangles. That is the folding intuition described by the source,
with congruence and coverage made explicit. Since \(\ell>0\), the
inner area is positive, and the ratio of outer to inner area is exactly two.
\end{proof}
\end{annexitem}
